\documentclass[border=10pt]{standalone}
\usepackage[siunitx, RPvoltages]{circuitikz}
\usepackage{tikz}

\begin{document}
\begin{circuitikz}[scale=0.8]
    % Input
    \draw (0,0) node[left]{Audio Input} to[short] (1,0);
    \draw (1,0) to[C=$C_1$] (3,0);
    
    % Transistor
    \draw (3,0) to[short] (4,0);
    \draw (4,-2) node[npn](npn){} node[below]{Q1};
    
    % Biasing resistors
    \draw (4,2) to[R, l_=$R_1$] (4,0);
    \draw (4,-2) to[R, l_=$R_E$] (4,-4);
    \draw (4,-4) to[short] (0,-4) to[battery1, l=9V] (0,2) to[short] (4,2);
    
    % Coupling capacitor
    \draw (5,0) to[C, l=$C_2$] (7,0);
    
    % Output
    \draw (7,0) to[short] (8,0) node[right]{Audio Output};
    
    % Bypass capacitor
    \draw (4,-2) to[C, l=$C_E$] (6,-2) to[short] (6,-4);
    
    % Labels
    \node[text width=3cm, align=center] at (4,-5.5) {\small Basic Audio Amplifier Circuit};
    \node[text width=4cm, align=center] at (8,-5.5) {\small For Improved Sound in Older TVs/Radios};
\end{circuitikz}
\end{document}